\documentclass[11pt]{article}
\usepackage[margin=25mm]{geometry}
\usepackage{fontspec,unicode-math}
\setmainfont{Latin Modern Roman}
\setmathfont{Latin Modern Math}
\usepackage{amsmath,amssymb,mathtools}
\usepackage{xurl,hyperref}
\hypersetup{colorlinks=true,urlcolor=blue,linkcolor=blue}
\setcounter{secnumdepth}{0}
\setlength{\parindent}{0pt}
\setlength{\parskip}{5pt}
\setlength{\emergencystretch}{3em}
\newcommand{\sref}[2]{\hyperlink{source:#1:#2}{[S#2]}}
\newcommand{\cataloguescope}[1]{\paragraph{Source question.}#1}
\newcommand{\floor}[1]{\left\lfloor#1\right\rfloor}
\begin{document}
\section{Degree-windowed Cech computation and the role of one b-function}\label{statement}

\subsection{Definitions and mathematical statement}

\paragraph{Shared notation.}Standard number systems are denoted by $\mathbb N,\mathbb Z,\mathbb Q,\mathbb R,\mathbb C$; index ranges, computational models and specialized notation follow the cited problem source and the definitions below.


For an ideal $I=(f_1,\ldots,f_r)$, the Cech complex uses localizations at products of its generators to compute local cohomology. Bernstein--Sato polynomials encode functional identities controlling powers and differential operators. The source asks whether fewer Cech computations or a single appropriate b-function computation can reduce the cost.
\paragraph{Statement recorded in the supplied source.}
Computing local cohomology via D-modules

The current method for computing local cohomology requires the iterated computation of Bernstein--Sato polynomials at every step of the Cech complex.

Can computing just the necessary steps of the Cech complex speed up the computation of local cohomology? What about computing just the Budur--Mustata--Saito b-function just once? \sref{statement}{2}
\paragraph{Scope qualification.}
The catalogue flags this record as reconstructed and unverified. The supplied question is preserved; the definitions do not certify a unique intended model. No universal one-b-function sufficiency claim is stated; the source names candidate computational improvements rather than supplying their hypotheses.
The supplied catalogue marks this formulation as an unverified reconstruction; the source link does not establish that all missing assumptions were recovered.
\paragraph{Source status.}
The supplied catalogue labels this record partially solved. Its source-status metadata was checked against the supplied snapshot on 2026-10-11. This import does not independently establish current open status or mathematical correctness.
\subsection{Short English statement}
Can local-cohomology computations avoid repeated b-function work by using only necessary Cech steps?
\subsection{Sources}
\begin{itemize}
\item[\hypertarget{source:statement:1}{S1}] Ulam AI and the UnsolvedMath Contributors, UnsolvedMath: A Curated Collection of Open Mathematics Problems, catalogue entry. Catalogue curation and metadata: CC BY 4.0. The supplied source question is retained; definitions and the English summary are editorial additions.
\url{https://huggingface.co/datasets/ulamai/UnsolvedMath}.
\item[\hypertarget{source:statement:2}{S2}] AIM Workshop Problem Lists, original problem source linked by the catalogue
\url{http://aimpl.org/macaulay2efie/2/}.
\end{itemize}
\paragraph{Preparation and review.}
Original question attributed to the cited source and UnsolvedMath contributors. Codex (GPT-6; exact serving revision not exposed) assisted with definitions, English summaries, source inspection and formatting. No new human mathematical authorship or independent proof verification is asserted. Catalogue identity checked against the complete supplied registry and existing repository statements; known flagged repetitions and obvious identical conjectures excluded. Full mathematical-equivalence/literature review has not been performed. Statement preparation and final review followed contribution prompts C and D; automated checks verify structure and identity, not open status.
\label{end:statement}
\end{document}
